\documentclass[aspectratio=169, 10pt]{beamer}

% --- Configuración de Tema y Colores ---
\usetheme{Singapore}
\usecolortheme{whale}
\setbeamertemplate{navigation symbols}{}

% --- Paquetes requeridos ---
\usepackage[utf8]{inputenc}
\usepackage[spanish, es-tabla]{babel}
\usepackage{listings}
\usepackage{xcolor}
\usepackage{tikz}
\usepackage{booktabs}

% --- Configuración para Listados de Código ---
\definecolor{codegreen}{rgb}{0,0.6,0}
\definecolor{codegray}{rgb}{0.5,0.5,0.5}
\definecolor{codepurple}{rgb}{0.58,0,0.82}
\definecolor{backcolour}{rgb}{0.96,0.96,0.97}

\lstdefinestyle{mystyle}{
    backgroundcolor=\color{backcolour},   
    commentstyle=\color{codegreen},
    keywordstyle=\color{blue}\bfseries,
    numberstyle=\tiny\color{codegray},
    stringstyle=\color{codepurple},
    basicstyle=\ttfamily\scriptsize,
    breakatwhitespace=false,         
    breaklines=true,                 
    captionpos=b,                    
    keepspaces=true,                 
    numbers=left,                    
    numbersep=5pt,                  
    showspaces=false,                
    showstringspaces=false,
    showtabs=false,                  
    tabsize=2,
    frame=lines,
    rulecolor=\color{black!20}
}
\lstset{style=mystyle}

% --- Información de la Portada ---
\title[Administrador de Música]{Administrador de Música}
\institute[ESFM - IPN]{Escuela Superior de Física y Matemáticas}
\date{\today}

\begin{document}

% ----------------------------------------------------
% 1. TÍTULO
% ----------------------------------------------------
\begin{frame}
    \titlepage
\end{frame}

% ----------------------------------------------------
% 2. ÍNDICE
% ----------------------------------------------------
\begin{frame}{Índice de Contenidos}
    \tableofcontents
\end{frame}

% ----------------------------------------------------
% 3. MODELADO DE BASE DE DATOS
% ----------------------------------------------------
\section{SQL}
\subsection{Modelado de Base de Datos}
    \begin{frame}{¿Cuál es el objetivo?}
        El objetivo de esta estructura es relacionar las entidades para normalizar la base de datos y optimizar el almacenamiento.\bigskip

        La tabla \textbf{albumes} funciona como la entidad principal del catálogo. Para evitar duplicidad de información, los registros no se insertan en un único bloque; primero se registran las entidades independientes y posteriormente se vinculan a través de tablas relacionales intermedias.\bigskip

        Gracias a esto, un artista puede tener varios álbumes, un álbum puede reunir varias canciones y una canción puede pertenecer a más de un lanzamiento, evitando en todo momento la redundancia de datos.
    \end{frame}% =====================================================
%                      ENTIDADES
% =====================================================
\subsection{Entidades/Tablas principales}
\begin{frame}[t]{Entidades}
    \centering
    \footnotesize
    \setlength{\tabcolsep}{3pt}
    \renewcommand{\arraystretch}{0.95}

    \begin{columns}[T]
        % ---------------- COLUMNA 1 ----------------
        \begin{column}{0.32\textwidth}
            \centering
            \textbf{\texttt{usuarios}}\par\vspace{0.05cm}
            \begin{tabular}{ll}
                \toprule
                \textbf{Campo} & \textbf{Tipo} \\
                \midrule
                \texttt{id\_usuario} & INT (PK) \\
                \texttt{username} & VARCHAR(50) \\
                \texttt{password\_hash} & VARCHAR(255) \\
                \texttt{created\_at} & TIMESTAMP \\
                \bottomrule
            \end{tabular}

            \vspace{0.25cm}
            \textbf{\texttt{artistas}}\par\vspace{0.05cm}
            \begin{tabular}{ll}
                \toprule
                \textbf{Campo} & \textbf{Tipo} \\
                \midrule
                \texttt{id\_artista} & INT (PK) \\
                \texttt{nombre} & VARCHAR(150) \\
                \bottomrule
            \end{tabular}
        \end{column}

        % ---------------- COLUMNA 2 ----------------
        \begin{column}{0.32\textwidth}
            \centering
            \textbf{\texttt{albumes}}\par\vspace{0.05cm}
            \begin{tabular}{ll}
                \toprule
                \textbf{Campo} & \textbf{Tipo} \\
                \midrule
                \texttt{id\_album} & INT (PK) \\
                \texttt{titulo} & VARCHAR(150) \\
                \texttt{descripcion} & TEXT \\
                \texttt{anio\_lanzamiento} & YEAR \\
                \texttt{portada\_ruta} & VARCHAR(255) \\
                \bottomrule
            \end{tabular}

            \vspace{0.25cm}
            \textbf{\texttt{notas}}\par\vspace{0.05cm}
            \begin{tabular}{ll}
                \toprule
                \textbf{Campo} & \textbf{Tipo} \\
                \midrule
                \texttt{id\_nota} & INT (PK) \\
                \texttt{contenido} & TEXT \\
                \texttt{fecha} & TIMESTAMP \\
                \bottomrule
            \end{tabular}
        \end{column}

        % ---------------- COLUMNA 3 ----------------
        \begin{column}{0.32\textwidth}
            \centering
            \textbf{\texttt{canciones}}\par\vspace{0.05cm}
            \begin{tabular}{ll}
                \toprule
                \textbf{Campo} & \textbf{Tipo} \\
                \midrule
                \texttt{id\_cancion} & INT (PK) \\
                \texttt{titulo} & VARCHAR(150) \\
                \texttt{duracion\_segundos} & INT \\
                \texttt{ruta\_archivo} & VARCHAR(255) \\
                \bottomrule
            \end{tabular}

            \vspace{0.25cm}
            \textbf{\texttt{playlists}}\par\vspace{0.05cm}
            \begin{tabular}{ll}
                \toprule
                \textbf{Campo} & \textbf{Tipo} \\
                \midrule
                \texttt{id\_playlist} & INT (PK) \\
                \texttt{nombre} & VARCHAR(100) \\
                \texttt{descripcion} & TEXT \\
                \texttt{created\_at} & TIMESTAMP \\
                \bottomrule
            \end{tabular}
        \end{column}
    \end{columns}
    \vspace{0.5cm}
    \begin{enumerate}
        \item TIMESTAMP: es un tipo de dato que registra el instante exacto (fecha y hora) en que ocurre un evento, como la creación o modificación de un registro
        \item PK: Llave primaria (Primary Key)
    \end{enumerate}
\end{frame}


% =====================================================
%                      RELACIONES
% =====================================================
\subsection{Relaciones}
\begin{frame}[t]{Relaciones}
    \centering
    \footnotesize
    \setlength{\tabcolsep}{3pt}
    \renewcommand{\arraystretch}{0.95}

    \begin{columns}[T]
        % ---------------- COLUMNA 1 ----------------
        \begin{column}{0.32\textwidth}
            \centering
            \textbf{\texttt{rel\_tracks\_album}}\par\vspace{0.05cm}
            \begin{tabular}{ll}
                \toprule
                \textbf{Campo} & \textbf{Tipo} \\
                \midrule
                \texttt{id\_album} & INT (PK) \\
                \texttt{id\_cancion} & INT (PK) \\
                \texttt{numero\_pista} & INT \\
                \bottomrule
            \end{tabular}

            \vspace{0.25cm}
            \textbf{\texttt{rel\_art\_track}}\par\vspace{0.05cm}
            \begin{tabular}{ll}
                \toprule
                \textbf{Campo} & \textbf{Tipo} \\
                \midrule
                \texttt{id\_artista} & INT (PK) \\
                \texttt{id\_cancion} & INT (PK) \\
                \texttt{rol} & VARCHAR(50) \\
                \bottomrule
            \end{tabular}
        \end{column}

        % ---------------- COLUMNA 2 ----------------
        \begin{column}{0.32\textwidth}
            \centering
            \textbf{\texttt{rel\_tracks\_playlist}}\par\vspace{0.05cm}
            \begin{tabular}{ll}
                \toprule
                \textbf{Campo} & \textbf{Tipo} \\
                \midrule
                \texttt{id\_playlist} & INT (PK) \\
                \texttt{id\_cancion} & INT (PK) \\
                \texttt{posicion} & INT \\
                \texttt{agregado\_en} & TIMESTAMP \\
                \bottomrule
            \end{tabular}

            \vspace{0.25cm}
            \textbf{\texttt{rel\_art\_album}}\par\vspace{0.05cm}
            \begin{tabular}{ll}
                \toprule
                \textbf{Campo} & \textbf{Tipo} \\
                \midrule
                \texttt{id\_artista} & INT (PK) \\
                \texttt{id\_album} & INT (PK) \\
                \bottomrule
            \end{tabular}
        \end{column}

        % ---------------- COLUMNA 3 ----------------
        \begin{column}{0.32\textwidth}
            \centering
            \textbf{\texttt{rel\_user\_playlist}}\par\vspace{0.05cm}
            \begin{tabular}{ll}
                \toprule
                \textbf{Campo} & \textbf{Tipo} \\
                \midrule
                \texttt{id\_usuario} & INT (PK) \\
                \texttt{id\_playlist} & INT (PK) \\
                \texttt{es\_propietario} & BOOLEAN \\
                \bottomrule
            \end{tabular}

            \vspace{0.25cm}
            \textbf{\texttt{rel\_notes\_album}}\par\vspace{0.05cm}
            \begin{tabular}{ll}
                \toprule
                \textbf{Campo} & \textbf{Tipo} \\
                \midrule
                \texttt{id\_nota} & INT (PK) \\
                \texttt{id\_album} & INT (PK) \\
                \bottomrule
            \end{tabular}
        \end{column}
    \end{columns}
    \begin{enumerate}
        \item BOOLEAN: es un tipo de dato que solo puede almacenar dos valores lógicos: verdadero (TRUE) o falso (FALSE)
    \end{enumerate}
\end{frame}
\subsection{Modelo de relación}
% =====================================================
%             EXPLICACIÓN DE LAS RELACIONES
% =====================================================
\begin{frame}[t]{Explicación de las Relaciones}
    \small
    \begin{itemize}
        \setlength{\itemsep}{6pt}

        \item \textbf{\texttt{rel\_tracks\_album}}: relaciona el
        \texttt{id\_cancion} con el \texttt{id\_album}, indicando
        además el número de pista (\texttt{numero\_pista}) de la
        canción dentro del álbum.

        \item \textbf{\texttt{rel\_tracks\_playlist}}: relaciona el
        \texttt{id\_cancion} con el \texttt{id\_playlist}, guardando
        su \texttt{posicion} en la lista y cuándo se agregó
        (\texttt{agregado\_en}).

        \item \textbf{\texttt{rel\_art\_track}}: relaciona el
        \texttt{id\_artista} con el \texttt{id\_cancion}, y el campo
        \texttt{rol} indica su participación (por ejemplo, intérprete
        o compositor).

        \item \textbf{\texttt{rel\_art\_album}}: relaciona el
        \texttt{id\_artista} con el \texttt{id\_album}, es decir,
        qué artistas pertenecen a cada álbum.

        \item \textbf{\texttt{rel\_notes\_album}}: relaciona el
        \texttt{id\_nota} con el \texttt{id\_album}, permitiendo
        asociar notas a un álbum.

        \item \textbf{\texttt{rel\_user\_playlist}}: relaciona el
        \texttt{id\_usuario} con el \texttt{id\_playlist}, y
        \texttt{es\_propietario} indica si el usuario es dueño de la
        playlist.
    \end{itemize}
\end{frame}
\subsection{Ejemplo}
% =====================================================
%                 EJEMPLO DE REGISTROS
% =====================================================
% =====================================================
%                 EJEMPLO DE REGISTROS
% =====================================================
\begin{frame}[t]{Ejemplo: ¿cómo se relacionan los registros?}
    \centering
    \footnotesize
    \setlength{\tabcolsep}{4pt}
    \renewcommand{\arraystretch}{1.0}

    % ---------- FILA 1: ENTIDADES ----------
    \begin{columns}[T]
        \begin{column}{0.22\textwidth}
            \centering
            \textbf{\texttt{artistas}}\par\vspace{0.1cm}
            \begin{tabular}{ll}
                \toprule
                \textbf{id\_artista} & \textbf{nombre} \\
                \midrule
                1 & Daft Punk \\
                \bottomrule
            \end{tabular}
        \end{column}

        \begin{column}{0.36\textwidth}
            \centering
            \textbf{\texttt{albumes}}\par\vspace{0.1cm}
            \begin{tabular}{lll}
                \toprule
                \textbf{id\_album} & \textbf{titulo} & \textbf{anio} \\
                \midrule
                10 & Discovery & 2001 \\
                \bottomrule
            \end{tabular}
        \end{column}

        \begin{column}{0.40\textwidth}
            \centering
            \textbf{\texttt{canciones}}\par\vspace{0.1cm}
            \begin{tabular}{lll}
                \toprule
                \textbf{id\_cancion} & \textbf{titulo} & \textbf{duracion} \\
                \midrule
                100 & One More Time & 320 \\
                101 & Digital Love  & 298 \\
                \bottomrule
            \end{tabular}
        \end{column}
    \end{columns}

    \vspace{0.4cm}
    {\small $\Downarrow$ \textit{las tablas de relación los conectan con el álbum} $\Downarrow$}
    \vspace{0.4cm}

    % ---------- FILA 2: RELACIONES ----------
    \begin{columns}[T]
        \begin{column}{0.36\textwidth}
            \centering
            \textbf{\texttt{rel\_art\_album}}\par\vspace{0.1cm}
            \begin{tabular}{ll}
                \toprule
                \textbf{id\_artista} & \textbf{id\_album} \\
                \midrule
                1 & 10 \\
                \bottomrule
            \end{tabular}

            \vspace{0.15cm}
            \textit{Daft Punk $\rightarrow$ Discovery}
        \end{column}

        \begin{column}{0.56\textwidth}
            \centering
            \textbf{\texttt{rel\_tracks\_album}}\par\vspace{0.1cm}
            \begin{tabular}{lll}
                \toprule
                \textbf{id\_album} & \textbf{id\_cancion} & \textbf{numero\_pista} \\
                \midrule
                10 & 100 & 1 \\
                10 & 101 & 2 \\
                \bottomrule
            \end{tabular}

            \vspace{0.15cm}
            \textit{One More Time y Digital Love $\rightarrow$ Discovery}
        \end{column}
    \end{columns}
\end{frame}
% =====================================================
%            RESULTADO FINAL: UNIENDO LAS TABLAS
% =====================================================
\begin{frame}[t,fragile]{El álbum completo: uniendo tablas}
    \centering
    \footnotesize
    \setlength{\tabcolsep}{4pt}
    \renewcommand{\arraystretch}{1.0}

    \begin{block}{¿Está incompleta la información del álbum?}
        \footnotesize
        No: así \textbf{optimizamos espacio}. Cada dato se guarda una sola
        vez en su tabla, y al consultarlo \textbf{unimos columnas de
        distintas tablas} para reconstruir el registro completo.
    \end{block}

    \vspace{0.2cm}
    \textbf{Resultado final (después de la consulta SQL)}\par\vspace{0.1cm}
    \begin{tabular}{lllll}
        \toprule
        \textbf{album} & \textbf{artista} & \textbf{anio} &
        \textbf{pista} & \textbf{cancion} \\
        \midrule
        Discovery & Daft Punk & 2001 & 1 & One More Time \\
        Discovery & Daft Punk & 2001 & 2 & Digital Love  \\
        \bottomrule
    \end{tabular}

    \vspace{0.3cm}
    \begin{columns}[T]
        \begin{column}{0.52\textwidth}
            \textbf{Consulta SQL}
            \begin{lstlisting}[language=SQL,basicstyle=\ttfamily\tiny]
                SELECT al.titulo, ar.nombre,
                       al.anio_lanzamiento,
                       rt.numero_pista, c.titulo
                FROM albumes al
                JOIN rel_art_album ra
                     ON ra.id_album = al.id_album
                JOIN artistas ar
                     ON ar.id_artista = ra.id_artista
                JOIN rel_tracks_album rt
                     ON rt.id_album = al.id_album
                JOIN canciones c
                     ON c.id_cancion = rt.id_cancion
                WHERE al.id_album = 10;
            \end{lstlisting}
        \end{column}

        \begin{column}{0.44\textwidth}
            \textbf{En lenguaje natural}
            \begin{itemize}
                \setlength{\itemsep}{2pt}
                \item Toma las columnas \texttt{titulo}, \texttt{nombre},
                \texttt{anio\_lanzamiento}, \texttt{numero\_pista} y
                \texttt{titulo} de las tablas \texttt{albumes},
                \texttt{artistas}, \texttt{rel\_tracks\_album} y
                \texttt{canciones}.
                \item Únelas usando las tablas de relación
                (\texttt{rel\_art\_album} y \texttt{rel\_tracks\_album}).
                \item Donde el \texttt{id\_album} sea igual a \texttt{10}.
            \end{itemize}
        \end{column}
    \end{columns}
\end{frame}
\section{Python}
\begin{frame}{Integración con Python}
    \centering
    \vfill
    {\Huge \textbf{Código en Python}}\par\vspace{0.4cm}
    {\Huge Funciones explicadas}
    \vfill
\end{frame}
\subsection{Archivos}
\begin{frame}{Archivos}
El proyecto se dividio en 3 archivos \textbf{Python:}
\begin{enumerate}
    \item funciones.py \par
            Las funciones que constantemente usaremos en el programa
    \item conexion.py\par
            La conexión con la base de datos.
    \item main.py\par
            El codigo donde integramos todo.
\end{enumerate}
\end{frame}
\subsection{Recursos: Libreria mysql.connector}
\begin{frame}[t]{Librería \texttt{mysql.connector}}
    \small
    \centering
    \footnotesize
    \setlength{\tabcolsep}{4pt}
    \renewcommand{\arraystretch}{1.0}

    \begin{block}{¿Qué es?}
        \footnotesize
        Librería oficial de MySQL para que Python pueda
        conectarse a la base de datos, ejecutar consultas SQL y recibir los
        resultados.
    \end{block}
    

    \vspace{0.3cm}
    \textbf{Funciones que usamos:}
    \vspace{0.1cm}

    \begin{center}
        \footnotesize
        \renewcommand{\arraystretch}{1.15}
        \begin{tabular}{p{0.32\textwidth} p{0.60\textwidth}}
            \toprule
            \textbf{Función} & \textbf{Explicación} \\
            \midrule
            \texttt{mysql.connector.connect()} &
            Abre la conexión con la base de datos. La usamos dentro de
            \texttt{connect\_to\_database()} (archivo \texttt{conexion.py}). \\

            \texttt{cnx.cursor()} &
            Crea el cursor, el objeto que envía las instrucciones SQL. \\

            \texttt{cursor.execute(sql, valores)} &
            Ejecuta una consulta. Los \texttt{\%s} se sustituyen por los
            valores, lo que evita inyección SQL. \\

            \texttt{cursor.fetchone()} &
            Devuelve una sola fila (en \texttt{verify\_user}). \\

            \texttt{cursor.fetchall()} &
            Devuelve todas las filas (en las funciones \texttt{show\_*}). \\

            \texttt{cursor.lastrowid} &
            Entrega el ID del último registro insertado. (esto es un atributo, no una función)\\

            \texttt{cnx.commit()} &
            Guarda los cambios de \texttt{INSERT} de forma permanente. \\

            \texttt{cnx.rollback()} &
            Deshace los cambios si algo falla (en \texttt{add\_playlist}). \\

            \texttt{cursor.close()}, \texttt{cnx.close()} &
            Cierran el cursor y la conexión en el bloque \texttt{finally}. \\
            \bottomrule
        \end{tabular}
    \end{center}
\end{frame}
\end{document}